\section{Objetivos}
% En objetivos.tex (PDF aparte) la Figura de la arquitectura no existe.
\ifdefined\soloobjetivos
\newcommand{\figarq}{la arquitectura propuesta del sistema}
\else
\newcommand{\figarq}{la Figura~\ref{fig:arquitectura}}
\fi

\subsection{Objetivo General}
\textbf{Evaluar el efecto de un prototipo de control de acceso con códigos QR
	dinámicos sobre el tiempo de validación, la fila y las incidencias de
	identidad en el acceso peatonal piloto del TESCHA, para aportar evidencia que
	permita decidir si conviene ampliarlo a otros accesos.}

Se plantea porque hoy no existen datos que midan el problema en el TESCHA, y
una decisión sobre el sistema debe apoyarse en evidencia y no en suposiciones.
El tiempo de validación es el que transcurre desde que el estudiante llega al
punto de revisión hasta que lo cruza, con revisión visual o con lectura del
QR. El objetivo pide evaluar el efecto, no demostrar una mejora: si las
mediciones no muestran diferencia, ese también es un resultado válido.

\subsection{Objetivos Específicos}
Los ocho objetivos específicos se reparten en las dos líneas de trabajo. Cada
uno indica qué se hará, cómo, para qué y por qué se hace.

\subsubsection{Desarrollo Tecnológico}
\begin{itemize}
	\item[\textbf{OT1}] \textbf{Diseñar la arquitectura del sistema} —lector
	      ESP32-CAM, ESP32 del torniquete, servidor, base de datos relacional y
	      panel web— a partir de \figarq, con un diagrama de componentes y un
	      modelo entidad-relación normalizado, para tener una base común en la que
	      cada intento de ingreso quede registrado. \emph{Porque} el acceso actual
	      no deja registro digital de quién entra ni a qué hora, y el prototipo
	      necesita una estructura que lo registre.
	\item[\textbf{OT2}] \textbf{Implementar la generación y la validación de
		      códigos QR dinámicos firmados con HMAC}, con una contraseña de un solo
	      uso de vigencia corta, HTTPS entre lector, servidor y panel, y pruebas con
	      tokens vencidos, repetidos y alterados y con desfases de reloj, para
	      impedir que una captura de pantalla o un mensaje modificado permitan el
	      ingreso. \emph{Porque} un código estático puede copiarse y compartirse,
	      mientras que la firma con contraseña de un solo uso resiste el reenvío y
	      la modificación del mensaje \parencite{subpratatsavee2015hmac}.
	\item[\textbf{OT3}] \textbf{Construir el prototipo funcional} que integre
	      lector, torniquete, servidor y panel web, programando la lectura, la
	      consulta al servidor y los reportes de ingresos por franja horaria,
	      intentos rechazados y estudiantes sin ingresos consecutivos, para contar
	      con un producto utilizable en el piloto. \emph{Porque} sin un prototipo
	      no hay sistema que evaluar, y porque directivos y Control Escolar hoy no
	      disponen de datos de ingreso para tomar decisiones.
	\item[\textbf{OT4}] \textbf{Verificar el desempeño técnico del prototipo}
	      midiendo el tiempo de captura, decodificación y consulta por Wi-Fi y la
	      tasa de lecturas fallidas al variar la iluminación, el brillo del celular
	      y la distancia al lector, para descartar, antes del piloto, que el lector
	      sea un cuello de botella. \emph{Porque} no se encontró literatura
	      arbitrada que reporte el tiempo de lectura de un código QR en un
	      ESP32-CAM, de modo que debe medirse con el prototipo.
\end{itemize}

\subsubsection{Investigación y Análisis de Datos}
\begin{itemize}
	\item[\textbf{OI1}] \textbf{Analizar la situación actual del acceso} —tasa de
	      llegada, tiempo de validación, utilización del acceso, fila máxima e
	      incidencias de identidad— con la observación de cinco días
	      (Apéndice~C), la encuesta (Apéndice~A) y las entrevistas
	      (Apéndice~B), para establecer la línea base con la que se compara el
	      piloto. \emph{Porque} no existe ningún registro del flujo de acceso del
	      TESCHA y, sin una línea base, no puede atribuirse ningún cambio al
	      sistema.
	\item[\textbf{OI2}] \textbf{Estimar la viabilidad y la aceptación del
		      sistema entre los estudiantes} —celular propio, conexión, batería y
	      disposición a mostrar un QR— con la encuesta del Apéndice~A, aplicada a
	      una muestra calculada para proporciones en población finita, para saber
	      si la población puede y quiere usar la solución. \emph{Porque} el sistema
	      depende de que el estudiante use su propio celular y se desconoce si la
	      comunidad lo aceptaría.
	\item[\textbf{OI3}] \textbf{Interpretar las obligaciones de protección de
		      datos personales aplicables al TESCHA}, revisando la ley estatal
	      \parencite{edomex2017datos} y la ley general \parencite{lgpdppso2025} y
	      consultando a Control Escolar en la entrevista del Apéndice~B, para
	      traducirlas en requisitos del sistema. \emph{Porque} el TESCHA es sujeto
	      obligado de ambas leyes y el prototipo registra la hora de ingreso de
	      cada estudiante.
	\item[\textbf{OI4}] \textbf{Evaluar el efecto del sistema en los tres
		      indicadores} contrastando la línea base con el piloto, con un diseño de
	      preprueba y posprueba de un solo grupo, el intervalo de cinco minutos
	      como unidad de análisis emparejada por posición y día (40 pares) y una
	      prueba para muestras relacionadas, para responder la pregunta de
	      investigación y recomendar si el piloto se amplía. \emph{Porque} la
	      pregunta pide el efecto del sistema, y eso solo se conoce comparando un
	      antes y un después.
\end{itemize}

\subsection{Relación con la Pregunta de Investigación}
Los objetivos de desarrollo tecnológico producen la variable independiente de
la pregunta, el mecanismo de validación basado en QR dinámico. Los objetivos
de investigación y análisis de datos miden las variables dependientes —tiempo
de validación, longitud de la fila e incidencias de identidad— y aportan lo que
el diseño necesita para interpretarlas: la línea base, la aceptación de los
estudiantes y las condiciones legales. Ningún objetivo queda sin función: sin
los primeros no hay sistema que evaluar y sin los segundos no hay forma de
saber si el sistema sirvió.
